%! Piecewise-Defined Functions — Unit 2, Lesson 2.6 — Group Activity
\documentclass[10pt]{article}
\usepackage{atda-article}
\usepackage{atda-boxes}
\newcommand{\TallMath}[1]{$\displaystyle #1\rule[-1.4em]{0pt}{3.2em}$}
\setlength{\workrowsep}{6pt}

\begin{document}

\pageheader{Unit 2, Lesson 2.6}{Group Activity --- Which Piece Owns This Input?}
% NAMESTRIP: no \namedateperiod here. The student writes their name once, on the
% cover this component is stapled behind. See references/conventions.md.

\vspace{0.15in}

\begin{scenariobox}[The Scenario]{royal}
\small
Every question today is answered by reading a graph. \textbf{The rule for the whole activity: before
you use a rule or read a value, say out loud which piece owns the input.} And whenever a boundary
point is involved, \textbf{let the dots decide} --- filled means included, open means not.
\textbf{Everyone starts at Tier R.} When your whole group can do Tier R without help, move on.
\end{scenariobox}

\vspace{0.10in}

\boxguard[30]
\begin{tcolorbox}[colback=white, colframe=black!40, title=\textbf{Tier R --- Remediate},
    fonttitle=\bfseries, arc=2mm, left=3mm, right=3mm, top=2mm, bottom=2mm, breakable]
\small
A phone plan costs \$30 a month for anything up to $10$ gigabytes of data, and \$5 for each gigabyte
beyond $10$. The graph shows $C(g)$, the monthly cost in dollars when $g$ gigabytes are used. The plan
shuts data off at $20$ gigabytes.

\begin{center}
\begin{tikzpicture}
\begin{axis}[
    width=10.4cm, height=4.6cm,
    xlabel={$g$ (gigabytes used)}, ylabel={$C$ (monthly cost, \$)},
    xmin=0, xmax=20, ymin=0, ymax=90,
    xtick={0,2,4,6,8,10,12,14,16,18,20}, ytick={0,10,20,30,40,50,60,70,80},
    minor x tick num=1, minor y tick num=1,
    grid=both, grid style={linegray!45}, minor grid style={linegray!30},
    axis lines=left,
    tick label style={font=\tiny}, label style={font=\scriptsize}]
  \addplot[royal, very thick] coordinates {(0,30) (10,30) (20,80)};
\end{axis}
\end{tikzpicture}
\end{center}

\begin{enumerate}[leftmargin=1.6em, itemsep=4pt, topsep=2pt]
  \item Read the graph and fill in both rows. For the second row write \emph{flat} or \emph{climbing}
        to say which piece owns that input.
\smallskip

\renewcommand{\arraystretch}{1.3}
\begin{tabular}{@{}|l|c|c|c|c|c|@{}}
\hline
$g$ (gigabytes) & $4$ & $10$ & $14$ & $18$ & $20$ \\ \hline
$C(g)$ (\$) & \blank{1.2cm} & \blank{1.2cm} & \blank{1.2cm} & \blank{1.2cm} & \blank{1.2cm} \\ \hline
Which piece? & \blank{1.4cm} & \blank{1.4cm} & \blank{1.4cm} & \blank{1.4cm} & \blank{1.4cm} \\
\hline
\end{tabular}

  \item Give the \emph{boundary point} --- the input where one rule hands off to the next --- and say
        in one sentence what it means to a customer.

\writelines{2}

  \item Name the interval of inputs on which $C$ is constant and the interval on which $C$ is
        increasing. Then say whether the graph is continuous, and how you can tell from the picture.

\writelines{3}

  \item Give the domain and the range of $C$. Then give the absolute maximum and the absolute
        minimum, each with a \emph{value and a location}. Read the minimum's location carefully ---
        it is not a single point.

\writelines{3}
\end{enumerate}
\end{tcolorbox}

\vspace{0.10in}

\boxguard[30]
\begin{tcolorbox}[colback=white, colframe=black!40,
    title=\textbf{Tier A --- Approaching Proficiency},
    fonttitle=\bfseries, arc=2mm, left=3mm, right=3mm, top=2mm, bottom=2mm, breakable]
\small
Graph (a) shows $G(h)$, the cost in dollars of parking in a city garage for $h$ hours: \$4 for up to
and including $1$ hour, \$7 for over $1$ up to and including $3$, and \$12 for over $3$ up to and
including $6$. Graph (b) shows a function $f$ with no context attached. \textbf{In both graphs, filled
circles are part of the graph and open circles are not.}

\begin{center}
\begin{tabular}{@{}c@{\hspace{0.5cm}}c@{}}
\begin{tikzpicture}
\begin{axis}[
    width=6.5cm, height=4.4cm, title={\scriptsize (a) $G$, garage cost},
    xlabel={$h$ (hours parked)}, ylabel={$G$ (\$)},
    xmin=0, xmax=6.4, ymin=0, ymax=15,
    xtick={0,1,2,3,4,5,6}, ytick={0,2,4,6,8,10,12,14},
    grid=both, grid style={linegray!45},
    axis lines=left,
    tick label style={font=\tiny}, label style={font=\tiny}]
  \addplot[royal, very thick] coordinates {(0,4) (1,4)};
  \addplot[royal, very thick] coordinates {(1,7) (3,7)};
  \addplot[royal, very thick] coordinates {(3,12) (6,12)};
  \addplot[royal, only marks, mark=*, mark size=1.9pt] coordinates {(1,4) (3,7) (6,12)};
  \addplot[royal, only marks, mark=*, mark size=1.9pt, mark options={fill=white}]
    coordinates {(0,4) (1,7) (3,12)};
\end{axis}
\end{tikzpicture}
&
\begin{tikzpicture}
\begin{axis}[
    width=6.3cm, height=4.4cm, title={\scriptsize (b) $f$},
    xmin=-3, xmax=7, ymin=-3, ymax=7,
    xtick={-3,-2,-1,0,1,2,3,4,5,6,7}, ytick={-2,0,2,4,6},
    minor y tick num=1,
    grid=both, grid style={linegray!45}, minor grid style={linegray!30},
    axis lines=left, tick label style={font=\tiny}]
  \addplot[royal, very thick, domain=-2:2, samples=60]{x^2};
  \addplot[royal, very thick] coordinates {(2,6) (6,-2)};
  \addplot[royal, only marks, mark=*, mark size=1.9pt] coordinates {(-2,4) (2,4) (6,-2)};
  \addplot[royal, only marks, mark=*, mark size=1.9pt, mark options={fill=white}]
    coordinates {(2,6)};
\end{axis}
\end{tikzpicture}
\end{tabular}
\end{center}

\begin{enumerate}[leftmargin=1.6em, itemsep=4pt, topsep=2pt]
  \item \textbf{Graph (a).} Fill in the costs, then answer the boundary question underneath.
\smallskip

\renewcommand{\arraystretch}{1.3}
\begin{tabular}{@{}|l|c|c|c|c|c|@{}}
\hline
$h$ (hours) & $1$ & $2$ & $3$ & $4$ & $6$ \\ \hline
$G(h)$ (\$) & \blank{1.2cm} & \blank{1.2cm} & \blank{1.2cm} & \blank{1.2cm} & \blank{1.2cm} \\
\hline
\end{tabular}
\smallskip

A driver parks for \emph{exactly} $3$ hours and is charged \$12. Using the dots, say what they should
have been charged and which circle proves it.

\writelines{2}

  \item \textbf{Graph (a).} Give the range of $G$. Then say whether $G$ is continuous, naming every
        input where you would have to lift your pencil.

\writelines{2}

  \item \textbf{Graph (a), the one that sounds wrong at first.} The cost clearly goes up as you park
        longer. Is there \emph{any} interval on which $G$ is increasing? Answer, and explain using
        what the graph does on each piece.

\writelines{3}

  \item \textbf{Graph (b).} Give $f(-2)$, $f(0)$, $f(2)$, $f(4)$ and $f(6)$. Then give the domain of
        $f$, its zeros, and its $y$-intercept.

\writelines{3}

  \item \textbf{Graph (b), the stretch.} Give the range of $f$, then give the absolute maximum and
        the absolute minimum. \emph{Look at the open circle before you answer} --- one of these two
        does not exist, and Lesson 2.5 gave you the words for why.

\writelines{3}
\end{enumerate}
\end{tcolorbox}

\vspace{0.10in}

\boxguard[30]
\begin{tcolorbox}[colback=white, colframe=black!40, title=\textbf{Tier E --- Extension},
    fonttitle=\bfseries, arc=2mm, left=3mm, right=3mm, top=2mm, bottom=2mm, breakable]
\small

\begin{enumerate}[leftmargin=1.6em, itemsep=4pt, topsep=2pt]
  \item \textbf{Two claims, and each one misreads a piece.} Name the mistake and give the correct
        statement.
        \begin{enumerate}[label=(\alph*), leftmargin=1.6em, itemsep=3pt, topsep=3pt]
          \item About the shipping graph in your notes: \emph{``At $w=2$ the cost is both \$6 and
                \$10, so $S$ is not a function.''}

\writelines{2}

          \item About the garden centre's pay: \emph{``A $46$-hour week pays $46 \times \$12$ plus
                $6 \times \$18$, which is \$660.''}

\writelines{2}
        \end{enumerate}

  \item \textbf{Possible or impossible?} For each description, either name a graph from today that
        does it, or explain why nothing can.
        \begin{enumerate}[label=(\alph*), leftmargin=1.6em, itemsep=3pt, topsep=3pt]
          \item A piecewise function that is continuous everywhere and still has a sharp corner.

\writelines{1}

          \item A piecewise function whose graph gives two different outputs at one input.

\writelines{2}

          \item A piecewise function that is constant on every one of its pieces, yet whose outputs
                are larger at the right end than at the left.

\writelines{2}
        \end{enumerate}

  \item \textbf{Write the rules.} A pool charges a flat \$3 to enter and then \$2 for each hour after
        the second. The graph of $h(x)$, the total cost in dollars for $x$ hours, is below. Write the
        two rules \emph{and} the stretch of inputs each one owns, then check your second rule by
        computing the cost at $x=6$ and comparing it with the graph.

\begin{center}
\begin{tikzpicture}
\begin{axis}[
    width=8.4cm, height=3.6cm,
    xlabel={$x$ (hours)}, ylabel={$h$ (total cost, \$)},
    xmin=0, xmax=6, ymin=0, ymax=13,
    xtick={0,1,2,3,4,5,6}, ytick={0,3,5,7,9,11},
    grid=both, grid style={linegray!45},
    axis lines=left,
    tick label style={font=\tiny}, label style={font=\tiny}]
  \addplot[royal, very thick] coordinates {(0,3) (2,3) (6,11)};
\end{axis}
\end{tikzpicture}
\end{center}

\writelines{4}
\end{enumerate}
\end{tcolorbox}

\end{document}
